\section{Target Adaptation and Ablations}
\label{sec:supp-adaptation}

\subsection{Label-Free Target Adaptation}
\label{sec:adaptation}

The source-threshold audit shows that a threshold trained on A is too quiet on B. A target site that can observe unlabelled traces can still match its requested rate without waiting for outcome labels: it chooses the most aggressive threshold whose empirical alarm rate on a warm-up half of the target subjects is at most the budget, then deploys that threshold on the held-out half. Table~\ref{tab:adaptation} compares this rate-matched rule with source transfer and with an oracle that uses target labels.

\begin{table}[H]
\centering
\caption{Target adaptation on site B, seed 0. The target cohort is split into disjoint warm-up and test halves. ``Source'' applies the A-calibrated threshold; ``rate-matched'' uses no target labels and matches the requested rate on the warm-up half; ``oracle'' recalibrates on target labels. Sensitivity and rate are evaluated on the target test half.}
\label{tab:adaptation}
\small
\begin{tabular}{lcccccc}
\toprule
 & \multicolumn{2}{c}{Source} & \multicolumn{2}{c}{Rate-matched} & \multicolumn{2}{c}{Oracle} \\
\cmidrule(lr){2-3}\cmidrule(lr){4-5}\cmidrule(lr){6-7}
Budget & Sens. & Rate & Sens. & Rate & Sens. & Rate \\
\midrule
0.25 & $0.194$ & $0.094$ & $0.322$ & $0.237$ & $0.261$ & $0.171$ \\
0.50 & $0.261$ & $0.171$ & $0.414$ & $0.356$ & $0.414$ & $0.356$ \\
1.00 & $0.414$ & $0.356$ & $0.620$ & $0.950$ & $0.620$ & $0.950$ \\
2.00 & $0.620$ & $0.950$ & $0.809$ & $1.792$ & $0.809$ & $1.792$ \\
\bottomrule
\end{tabular}
\end{table}

On the target test half the label-free rule recovers the oracle operating point at 0.5, 1 and 2 per day: its threshold coincides with the target-calibrated threshold on the fixed grid. At 1 per day it raises sensitivity on that half from $0.414$ to $0.620$ and the realised rate from $0.356$ to $0.950$, still below the request. The source threshold is quieter on this half than on the full B test set, where the external audit reports $0.315$ sensitivity and a $0.276$ rate, because the warm-up/test partition changes the eligible-event mix. At 0.25 per day it selects a more aggressive threshold than the certified oracle, raising sensitivity from $0.194$ to $0.322$ at a rate of $0.237$ rather than the oracle's $0.261$ at $0.171$. This is a positive engineering result, but it must be read narrowly. The rule is transductive: it uses a warm-up sample from the deployment site, and it controls the alarm rate rather than the risk. If the target case mix changes the relationship between score and event risk, a rate-matched threshold can meet the budget while detecting fewer events. The oracle agreement here is a property of this cohort and this grid, not a general guarantee. The finite-sample certificate still requires labelled target calibration patients; rate matching is a deployment heuristic that can be audited with the operational metrics above.

\subsection{Ablations}
\label{sec:ablation}

One-at-a-time ablations on the primary A split identify which design choices are first-order. The refractory window is one: moving from $\tau=2$ to $\tau=12$ raises sensitivity from $0.518$ to $0.726$ but also raises the patient-violation fraction from $0.167$ to $0.327$ at the 1-per-day budget. The escalation tier is a strong de-alarming control: $\kappa=1$ lowers the realised rate from $0.631$ to $0.472$ and sensitivity from $0.607$ to $0.286$. The lead window changes sensitivity modestly ($0.560$ at 6 hours, $0.607$ at 12, $0.631$ at 24) without changing the threshold. The risk level is second-order on this split: $\alpha=0.05,0.1,0.2$ selects the same threshold, with certified risk changing from $0.844$ to $0.823$. Grid resolution is first-order at the coarse end: 20 thresholds produce a rate of $0.082$ and sensitivity $0.179$, 50 and 100 thresholds produce about $0.44$ and $0.51$, and 200 thresholds recover $0.631$ and $0.607$. The minimum-stay restriction changes both eligibility and the selected threshold: at 6, 12 and 48 hours the sensitivities are $0.627$, $0.655$ and $0.609$ at rates $0.832$, $0.825$ and $0.747$. These results do not justify a universal tuning rule, but they show that the reported operating point depends on declared structural choices, which is why those choices are fixed before calibration and reported with the certificate.

\subsection{Human-Factors Queueing Sensitivity}
\label{sec:queue}

The hourly capacity model above is a deliberately conservative staffing abstraction: one alarm per hour is equivalent to a 60-minute service time, far longer than a typical bedside response. To ground the translation in human-factors data, we add a continuous-time queue in which alarms wait for one of $c$ nurses and service times are exponential with a declared mean. Bridi et al.\ report a median alarm response of about 5 minutes in an intensive-care unit, with more than 60\% of alarms unanswered after 10 minutes during busy periods \citep{bridi2014reaction}. We therefore evaluate mean service times of 5 and 15 minutes; the latter represents detection, travel and bedside confirmation rather than the alarm acknowledgement alone.

\begin{table}[H]
\centering
\caption{Alarm-handling queue on a simulated 20-bed ward with five admissions per day. Service times are exponential; ``Util.'' is nurse utilisation over the 30-day horizon. The 5-minute mean is the median response reported by Bridi et al.; 15 minutes adds travel and confirmation. At 0.25 per day at most 0.7\% of CRC alarms wait more than 5 minutes.}
\label{tab:queue}
\small
\begin{tabular}{lcccccc}
\toprule
Budget & Policy & Service & Util. & $>5$ min & $>10$ min & p95 wait \\
\midrule
1.00 & Blocked & 5 min & $0.032$ & $0.000$ & $0.000$ & $0.0$ \\
1.00 & Blocked & 15 min & $0.095$ & $0.020$ & $0.015$ & $0.0$ \\
1.00 & CRC & 5 min & $0.050$ & $0.022$ & $0.006$ & $0.0$ \\
1.00 & CRC & 15 min & $0.150$ & $0.089$ & $0.070$ & $18.7$ \\
2.00 & Both & 5 min & $0.057$ & $0.024$ & $0.005$ & $0.2$ \\
2.00 & Both & 15 min & $0.171$ & $0.126$ & $0.094$ & $26.8$ \\
\bottomrule
\end{tabular}
\end{table}

Table~\ref{tab:queue} shows that the queue is not saturated even at the loosest budget, but waiting is not negligible once service includes travel. At 1 alarm per patient-day, the blocked policy keeps the fraction waiting more than 10 minutes below 2\% with a 15-minute service mean, whereas CRC reaches 7.0\% and a 95th-percentile wait of 18.7 minutes. At 2 per day both policies produce the same traces on this split and 9.4\% of alarms wait more than 10 minutes with one nurse; a second nurse removes essentially all waits. These fractions are lower than the 60\% unanswered rate in the observational study because our queue assumes perfect work-conserving service and no competing nursing tasks. The comparison is therefore a lower bound on human delay, not a validation of the ward model. Its engineering value is to show that the difference between a 1-per-day and a 2-per-day budget is visible in response latency even when average nurse utilisation is below 20\%, because alarm bursts matter more than the daily mean.
